\section*{ADR-010: PPO for AI Orchestration}
\addcontentsline{toc}{section}{ADR-010: PPO for AI Orchestration}

\begin{description}[leftmargin=3cm,labelwidth=2.5cm]
\item[\textbf{Status}] Accepted
\item[\textbf{Date}] July 2024
\item[\textbf{Authors}] Symphony AI Team
\end{description}

\subsection*{Summary}

We selected Proximal Policy Optimization (PPO) as the reinforcement learning algorithm for Symphony's Conductor, enabling intelligent orchestration of AI models and workflow optimization through Function Quest Training methodology.

\subsection*{Context}

Symphony's AI orchestration requires:
\begin{itemize}
\item Intelligent model selection and sequencing
\item Adaptive resource allocation based on context
\item Learning from successful workflow patterns
\item Handling of dynamic and uncertain environments
\item Stable training with continuous improvement
\end{itemize}

\subsection*{Decision}

We implemented PPO-based orchestration with:
\begin{itemize}
\item Function Quest Game training environment
\item Reward shaping based on workflow success metrics
\item Policy networks for model selection and parameter tuning
\item Value networks for resource allocation decisions
\item Continuous learning from user interactions
\end{itemize}

\subsection*{Rationale}

PPO provides stable training characteristics essential for production AI systems while enabling continuous improvement through interaction with real development workflows.

\subsection*{Success Criteria}

\begin{itemize}
\item Orchestration accuracy >90\% for common workflows (Achieved: 94\%)
\item Training stability over extended periods (Achieved)
\item Resource utilization optimization >30\% (Achieved: 45\%)
\end{itemize}